\section{Introduction}\label{sec:intro}

Genes act through their interaction partners, and a tumour can differ from the tissue around it not only in which genes change but in how the changed genes sit among their partners. Protein interaction networks such as STRING \citep{szklarczyk2023} make this arrangement measurable: the change of each gene between a tumour and its matched normal tissue becomes a function on the vertices of a fixed graph, and the question becomes whether that function has a reproducible shape on the graph.

Persistent homology \citep{edelsbrunner2002} is a natural language for such questions. Given a function on a simplicial complex, a lower-star filtration adds simplices in the order of their function values, and the persistence diagram records when connected components ($\Hzero$) and loops ($\Hone$) appear and disappear. Diagrams are stable under perturbation of the function \citep{cohensteiner2007}, and they summarise every threshold at once instead of requiring one cut-off. Persistent homology has been applied to cancer-specific regulatory networks \citep{masoomy2021}, to the identification of key genes in cancer networks \citep{ramos2025}, to protein interaction networks in complex disease \citep{song2023}, to pathway networks combined with expression \citep{abdullahi2025}, to the topology of molecular landscapes for therapy selection \citep{rietman2017,benzekry2015}, and, in a related form, to brain networks \citep{lee2012}. A persistent loop is then often interpreted as a biologically meaningful structure, for example a ring of interacting genes around a protein that behaves differently.

Two questions are seldom asked in this literature. First, are the tests calibrated? Many applications compare diagrams with ad hoc permutation schemes, or with no null model at all, and the unit of inference (the patient or the gene) is not always explicit. Second, does $\Hone$ detect anything that a simpler statistic would not? A loop on a network is formed by vertices that enter the filtration before the vertices that close it. The same arrangement can arise from a connected module of changed genes, from changes concentrated at high-degree genes, or from spatial autocorrelation of the change on the graph, all of which can be measured without topology: by the size of connected components, by $\Hzero$ persistence, or by Moran's $I$ \citep{moran1950}. Protein networks also have strong degree heterogeneity and interconnected hubs \citep{maslov2002,colizza2006}, which can create or remove loops without any specific biology.

Here we address both questions with a paired design. For each patient we place the signed tumour-minus-normal change on a fixed STRING network and compute two lower-star filtrations of its clique complex, one admitting the most increased genes first and one the most decreased first. The difference between their total finite $\Hone$ persistence is a patient-level asymmetry that reverses sign exactly when the two specimens of a patient are exchanged. This gives an exact sign-flip test with the patient as the unit of inference. We attribute persistence to genes through exact cycle representatives, checked against the persistence pairing, and control the family-wise error rate over all genes. We define a network-adjusted score that compares each patient's asymmetry with relabellings of that patient's own changes within strata of network degree and clustering, and we prove that it keeps the exact sign reversal.

We then evaluate the framework honestly. We calibrate every test under the null hypothesis on subnetworks of the real STRING graph with four noise distributions, first on development seeds and then, with the final code, on reserved verification seeds. And we design a matched experiment in which two patterns of change share the same network region, the same number of changed genes, the same change sizes and the same noise, and differ only in whether the changed genes enclose a hole around an unchanged gene. Comparing $\Hone$ with $\Hzero$, component size, Moran's $I$ and a topology-free local-extremum statistic on this experiment shows what directional persistence on a network actually measures.

Our results are partly negative, and we report them as such. The tests are well calibrated. Directional $\Hone$ persistence does carry a hole-specific signal that grows with effect size and that $\Hzero$ and the local-extremum statistic lack, but the raw asymmetry is dominated by connected blocks of change, and our degree- and clustering-matched adjustment does not remove the effect of change concentrated at hubs. We argue that these facts matter for every application of persistent homology to functions on molecular networks, and we give the tests and the benchmark needed to check them. The framework is then applied to \cohortN{} TCGA breast tumour--normal pairs with RNA and promoter-methylation data \citep{tcga2012} under an analysis plan frozen before any outcome is computed. \pending{TCGA results.}

The contribution is methodological. We do not introduce a new topological construction; we combine standard lower-star persistence \citep{edelsbrunner2002,cohensteiner2007} with a paired, sign-equivariant design, explicit null hypotheses, an attribution step with family-wise error control, and a benchmark that separates topology from simpler network structure. To our knowledge this combination, and in particular the matched evaluation against non-topological baselines, has not been reported for molecular change on interaction networks.
